\documentclass[12pt,a4paper]{article}
\usepackage[utf8]{inputenc}
\usepackage[spanish]{babel}
\usepackage{geometry}
\usepackage{graphicx}
\usepackage{hyperref}
\usepackage{enumitem}
\usepackage{booktabs}
\usepackage{listings}
\usepackage{xcolor}

\geometry{margin=2.5cm}
\hypersetup{
    colorlinks=true,
    linkcolor=blue,
    urlcolor=blue,
    citecolor=blue
}

\title{\textbf{Especificación del Proyecto}\\
\large Secure Fair: Sistema de Registro Criptográficamente Controlado para Ferias de Servicio Social}
\author{Equipo de Desarrollo}
\date{Febrero 2026}

\begin{document}

\maketitle
\newpage

\tableofcontents
\newpage

\section{Contexto}

Cada semestre, el departamento de Servicio Social de la universidad organiza una feria donde los socioformadores presentan sus proyectos y los estudiantes se registran para participar. Actualmente, los estudiantes se registran para un horario específico para asistir a la feria (evitando aglomeraciones) y reciben un código QR para validar su entrada. Cada proyecto muestra un código QR estático que enlaza a un formulario de registro.

Sin embargo, el sistema actual presenta desafíos operativos y de seguridad que requieren una solución integral y técnicamente robusta.

\section{Planteamiento del Problema}

El proceso actual sufre de cuatro problemas principales:

\subsection{Compartir Códigos QR}

Los estudiantes pueden tomar fotografías de los códigos QR de los proyectos y enviarlos a otros estudiantes, permitiendo registros remotos sin asistir físicamente a la feria. Esto elimina el propósito de la interacción presencial entre estudiantes y socioformadores.

\subsection{Falta de Verificación de Presencia Física}

Aunque los estudiantes deben reservar un horario, el proceso de registro al proyecto no está vinculado a la asistencia física verificada en la feria.

\subsection{Registros Múltiples}

Algunos estudiantes se registran en múltiples proyectos, lo cual viola las políticas de Servicio Social que requieren un único proyecto por periodo.

\subsection{Control de Capacidad}

Los proyectos con vacantes limitadas reciben aplicaciones excesivas sin mecanismos de control en tiempo real, utilizando formularios estáticos sin validación de capacidad o inscripción.

\section{Objetivos del Proyecto}

El sistema debe cumplir con los siguientes objetivos:

\begin{itemize}[leftmargin=*]
    \item Exigir asistencia física verificada antes del registro al proyecto
    \item Prevenir que los estudiantes se inscriban en múltiples proyectos dentro del mismo periodo
    \item Aplicar límites de capacidad por proyecto en tiempo real
    \item Proporcionar visibilidad operativa en tiempo real al departamento de Servicio Social
    \item Ser lo suficientemente simple para usuarios no técnicos (socioformadores)
    \item Incluir componentes criptográficos significativos como requisito del curso
\end{itemize}

\section{Alcance}

\subsection{Incluido en el MVP}

\begin{itemize}[leftmargin=*]
    \item Sistema de autenticación basado en roles (Administrador, Socioformador, Estudiante)
    \item Gestión de organizaciones, proyectos, periodos y franjas horarias
    \item Sistema de registro de estudiantes a franjas horarias con control de capacidad
    \item Generación de códigos QR firmados para entrada
    \item Verificación de entrada (check-in) con escaneo o entrada manual
    \item Generación de códigos de inscripción de un solo uso con caducidad
    \item Proceso de inscripción condicionado a verificación física
    \item Restricción de un proyecto por periodo por estudiante
    \item Panel de métricas y analíticas para administradores
    \item Exportación de datos maestros y por proyecto
    \item Implementación criptográfica (recibos de inscripción firmados digitalmente)
\end{itemize}

\subsection{Excluido del MVP}

\begin{itemize}[leftmargin=*]
    \item Lista de espera automática cuando los proyectos están llenos
    \item Importación de datos mediante Excel (se considerará para iteraciones posteriores)
    \item Sistema de notificaciones por correo o SMS
    \item Aplicación móvil nativa
    \item Integración con sistemas institucionales externos
\end{itemize}

\section{Roles del Sistema}

\subsection{Estudiante}

\textbf{Responsabilidades:}
\begin{itemize}[leftmargin=*]
    \item Registrarse en una franja horaria disponible
    \item Obtener código QR de entrada
    \item Verificar su entrada en la feria
    \item Ingresar código de inscripción proporcionado por el socioformador
    \item Aceptar las reglas del proyecto
    \item Recibir confirmación de inscripción firmada
\end{itemize}

\textbf{Restricciones:}
\begin{itemize}[leftmargin=*]
    \item Solo puede registrarse en un proyecto por periodo de feria
    \item Debe estar físicamente verificado (check-in) para inscribirse
    \item No puede compartir códigos de inscripción (son de un solo uso)
\end{itemize}

\subsection{Socioformador}

\textbf{Responsabilidades:}
\begin{itemize}[leftmargin=*]
    \item Acceder al panel con proyectos asignados
    \item Generar códigos de inscripción de corta duración
    \item Mostrar código en stand físico
    \item Ver conteo en tiempo real de inscripciones
    \item Exportar lista de estudiantes inscritos
\end{itemize}

\textbf{Características del diseño:}
\begin{itemize}[leftmargin=*]
    \item Interfaz minimalista diseñada para usuarios no técnicos
    \item Operación sin conocimientos de programación
    \item Generación de códigos con un clic
\end{itemize}

\subsection{Administrador (Servicio Social)}

\textbf{Responsabilidades:}
\begin{itemize}[leftmargin=*]
    \item Crear y gestionar periodos de feria
    \item Gestionar organizaciones y proyectos
    \item Definir franjas horarias con capacidades
    \item Asignar socioformadores a proyectos
    \item Monitorear asistencia en tiempo real
    \item Revisar analíticas de inscripciones
    \item Exportar tablas maestras de registro
    \item Realizar verificación de entrada (check-in)
\end{itemize}

\section{Arquitectura del Sistema}

\subsection{Stack Tecnológico}

\begin{itemize}[leftmargin=*]
    \item \textbf{Frontend:} React con Vite y TypeScript, Material-UI o shadcn/ui, React Router, React Query
    \item \textbf{Backend:} FastAPI (Python), SQLAlchemy 2.0, Alembic para migraciones
    \item \textbf{Base de Datos:} PostgreSQL con restricciones transaccionales
    \item \textbf{Autenticación:} JWT con tokens de acceso de corta duración
    \item \textbf{Criptografía:} Ed25519 para firmas digitales usando pynacl o cryptography
    \item \textbf{Hashing:} argon2-cffi para contraseñas, HMAC para códigos
    \item \textbf{DevOps:} Docker Compose, despliegue en Render/Fly.io/Railway
\end{itemize}

\subsection{Diagrama de Alto Nivel}

\begin{verbatim}
+-------------+         +-------------+         +--------------+
|   Cliente   | HTTPS   |   FastAPI   |  SQL    |  PostgreSQL  |
|   React     +-------->|   Backend   +-------->|   Database   |
|   (SPA)     |<--------+   (API)     |<--------+              |
+-------------+         +-------------+         +--------------+
                              |
                              | Crypto Operations
                              v
                        +-------------+
                        |  Ed25519    |
                        |  Signing    |
                        +-------------+
\end{verbatim}

\section{Lista de Características}

\subsection{Características Centrales (MVP)}

\begin{enumerate}[leftmargin=*]
    \item \textbf{Autenticación y Control de Acceso Basado en Roles (RBAC)}
    \begin{itemize}
        \item Login/logout con JWT
        \item Protección de endpoints por rol
        \item Verificación de permisos en servidor
    \end{itemize}
    
    \item \textbf{Gestión de Periodos y Franjas Horarias}
    \begin{itemize}
        \item CRUD de periodos de feria
        \item Definición de franjas horarias con capacidad
        \item Control de capacidad en tiempo real
    \end{itemize}
    
    \item \textbf{Gestión de Organizaciones y Proyectos}
    \begin{itemize}
        \item CRUD de organizaciones socioformadoras
        \item CRUD de proyectos con descripción y reglas
        \item Asignación de capacidad por proyecto
        \item Vinculación de socioformadores a proyectos
    \end{itemize}
    
    \item \textbf{Registro a Franjas Horarias}
    \begin{itemize}
        \item Estudiantes reservan franja horaria
        \item Validación de capacidad antes de confirmar
        \item Restricción de registro único por estudiante
    \end{itemize}
    
    \item \textbf{Sistema de Verificación de Entrada}
    \begin{itemize}
        \item Generación de QR firmado para cada reservación
        \item Endpoint de check-in para administradores
        \item Marca temporal de verificación
        \item Dashboard de asistencia en tiempo real
    \end{itemize}
    
    \item \textbf{Generación de Códigos de Inscripción}
    \begin{itemize}
        \item Socioformador genera código alfanumérico corto
        \item Almacenamiento como hash HMAC (no texto plano)
        \item Caducidad configurable (60-120 segundos)
        \item Un solo uso por código
    \end{itemize}
    
    \item \textbf{Inscripción Condicionada}
    \begin{itemize}
        \item Estudiante ingresa código manualmente
        \item Validación de check-in previo obligatorio
        \item Verificación de inscripción única por periodo
        \item Validación de capacidad disponible
        \item Transacción atómica para prevenir condiciones de carrera
    \end{itemize}
    
    \item \textbf{Recibos de Inscripción Firmados}
    \begin{itemize}
        \item Generación de recibo con firma Ed25519
        \item Almacenamiento de firma en base de datos
        \item Endpoint de verificación de integridad
    \end{itemize}
    
    \item \textbf{Analíticas y Reportes}
    \begin{itemize}
        \item Dashboard con métricas clave
        \item Gráficas de asistencia por franja horaria
        \item Inscripciones por proyecto
        \item Tasas de ocupación
    \end{itemize}
    
    \item \textbf{Exportación de Datos}
    \begin{itemize}
        \item Exportación maestra en CSV/XLSX
        \item Exportación por proyecto para socioformadores
        \item Filtros por periodo y organización
    \end{itemize}
\end{enumerate}

\subsection{Características Extendidas (Post-MVP)}

\begin{itemize}[leftmargin=*]
    \item Importación de proyectos mediante plantilla Excel
    \item Notificaciones por correo electrónico
\end{itemize}

\section{Flujos de Alto Nivel}

\subsection{Flujo de Registro a Franja Horaria}

\begin{enumerate}
    \item Estudiante ingresa al sistema
    \item Visualiza franjas horarias disponibles del periodo activo
    \item Selecciona franja con cupos disponibles
    \item Sistema valida capacidad en tiempo real
    \item Confirma registro
    \item Sistema genera código QR firmado
    \item Estudiante descarga o visualiza QR
\end{enumerate}

\subsection{Flujo de Verificación de Entrada}

\begin{enumerate}
    \item Estudiante llega a la entrada de la feria
    \item Muestra código QR al personal de Servicio Social
    \item Administrador escanea o ingresa token manualmente
    \item Sistema verifica firma digital y validez temporal
    \item Sistema marca estudiante como verificado (check-in)
    \item Estudiante recibe confirmación de acceso
\end{enumerate}

\subsection{Flujo de Inscripción a Proyecto}

\begin{enumerate}
    \item Socioformador genera código desde su panel
    \item Sistema crea código aleatorio, lo hashea y establece caducidad
    \item Código se muestra en pantalla del socioformador
    \item Estudiante (verificado previamente) ingresa código en su aplicación
    \item Sistema valida:
    \begin{itemize}
        \item Estudiante tiene check-in activo
        \item Código es válido y no ha expirado
        \item Código no ha sido usado
        \item Estudiante no está inscrito en otro proyecto del periodo
        \item Proyecto tiene capacidad disponible
    \end{itemize}
    \item Si todas las validaciones pasan:
    \begin{itemize}
        \item Se crea registro de inscripción
        \item Se marca código como usado
        \item Se genera recibo firmado digitalmente
        \item Se decrementa capacidad disponible
    \end{itemize}
    \item Estudiante recibe confirmación con recibo firmado
\end{enumerate}

\section{Modelo de Seguridad}

\subsection{Principios de Seguridad}

\begin{itemize}[leftmargin=*]
    \item \textbf{Defensa en profundidad:} Validación tanto en cliente como en servidor
    \item \textbf{Privilegio mínimo:} Cada rol solo accede a recursos necesarios
    \item \textbf{Fail-secure:} Fallos resultan en denegación de acceso, no en permisos
    \item \textbf{Separación de responsabilidades:} Lógica crítica aislada en capa de negocio
\end{itemize}

\subsection{Control de Acceso}

\textbf{Implementación RBAC:}
\begin{itemize}[leftmargin=*]
    \item Roles definidos en base de datos (ADMIN, SOCIO, STUDENT)
    \item Decoradores en endpoints verifican rol requerido
    \item Tokens JWT incluyen rol en payload
    \item Validación de pertenencia (socioformador solo ve sus proyectos)
\end{itemize}

\subsection{Protección de Endpoints}

\begin{itemize}[leftmargin=*]
    \item Autenticación requerida en todos los endpoints excepto /auth/login
    \item Validación de rol específico por endpoint
    \item Rate limiting en endpoints críticos (redención de códigos, check-in)
    \item Validación de entrada con pydantic schemas
\end{itemize}

\subsection{Protección contra Ataques Comunes}

\begin{itemize}[leftmargin=*]
    \item \textbf{Fuerza bruta:} Límite de intentos de login, códigos de inscripción
    \item \textbf{Inyección SQL:} Uso de ORM parametrizado (SQLAlchemy)
    \item \textbf{XSS:} Sanitización en frontend, Content-Security-Policy headers
    \item \textbf{CSRF:} Tokens JWT en headers (no cookies), SameSite si se usan cookies
    \item \textbf{Timing attacks:} Comparación constante de hashes
    \item \textbf{Condiciones de carrera:} Transacciones con SELECT FOR UPDATE
\end{itemize}

\section{Diseño Criptográfico}

\subsection{Hashing de Contraseñas}

\begin{itemize}[leftmargin=*]
    \item \textbf{Algoritmo:} Argon2id
    \item \textbf{Librería:} argon2-cffi
    \item \textbf{Parámetros:} Valores por defecto seguros
    \item \textbf{Propósito:} Protección de credenciales en base de datos
\end{itemize}

\subsection{Códigos de Inscripción}

\begin{itemize}[leftmargin=*]
    \item \textbf{Generación:} Caracteres alfanuméricos aleatorios (6-8 caracteres)
    \item \textbf{Almacenamiento:} HMAC-SHA256 con clave secreta del servidor
    \item \textbf{Verificación:} Comparación de HMAC en tiempo constante
    \item \textbf{Caducidad:} 60-120 segundos desde generación
    \item \textbf{Un solo uso:} Marcado como usado tras redención exitosa
\end{itemize}

\subsection{Recibos de Inscripción Firmados}

\begin{itemize}[leftmargin=*]
    \item \textbf{Algoritmo:} Ed25519 (firma de curva elíptica)
    \item \textbf{Librería:} pynacl o cryptography
    \item \textbf{Contenido del recibo:}
    \begin{verbatim}
    {
        "student_id": "<uuid>",
        "project_id": "<uuid>",
        "period_id": "<uuid>",
        "timestamp": "<ISO-8601>",
        "nonce": "<random>"
    }
    \end{verbatim}
    \item \textbf{Proceso:}
    \begin{enumerate}
        \item Servidor serializa datos del recibo (JSON canónico)
        \item Calcula firma Ed25519 con clave privada del servidor
        \item Almacena firma en base de datos
        \item Retorna recibo + firma al cliente
    \end{enumerate}
    \item \textbf{Verificación:}
    \begin{enumerate}
        \item Cliente o auditor envía recibo + firma
        \item Servidor re-serializa datos
        \item Verifica firma con clave pública
        \item Retorna válido/inválido
    \end{enumerate}
    \item \textbf{Propósito:} Integridad, autenticidad y no repudio
\end{itemize}

\subsection{Tokens QR de Entrada}

\begin{itemize}[leftmargin=*]
    \item \textbf{Contenido:} student\_id, slot\_id, timestamp, firma
    \item \textbf{Formato:} JWT firmado con clave del servidor o estructura similar
    \item \textbf{Caducidad:} Válido para el día del slot reservado
    \item \textbf{Verificación:} Validación de firma antes de check-in
\end{itemize}

\section{Modelo de Amenazas}

\subsection{Amenazas Identificadas}

\begin{enumerate}
    \item \textbf{Compartir QR de proyecto}
    \begin{itemize}
        \item \textit{Mitigación:} Códigos efímeros de un solo uso en lugar de QR estático
    \end{itemize}
    
    \item \textbf{Registro remoto sin asistencia física}
    \begin{itemize}
        \item \textit{Mitigación:} Gate de check-in obligatorio antes de inscripción
    \end{itemize}
    
    \item \textbf{Inscripción múltiple en proyectos}
    \begin{itemize}
        \item \textit{Mitigación:} Restricción UNIQUE en DB (student\_id, period\_id)
    \end{itemize}
    
    \item \textbf{Exceso de capacidad}
    \begin{itemize}
        \item \textit{Mitigación:} Transacciones atómicas con verificación de capacidad
    \end{itemize}
    
    \item \textbf{Fuerza bruta en códigos}
    \begin{itemize}
        \item \textit{Mitigación:} Rate limiting + caducidad + complejidad del código
    \end{itemize}
    
    \item \textbf{Suplantación de identidad}
    \begin{itemize}
        \item \textit{Mitigación:} Autenticación JWT + verificación en check-in
    \end{itemize}
    
    \item \textbf{Manipulación de recibos}
    \begin{itemize}
        \item \textit{Mitigación:} Firmas Ed25519 verificables
    \end{itemize}
    
    \item \textbf{Condiciones de carrera en capacidad}
    \begin{itemize}
        \item \textit{Mitigación:} SELECT FOR UPDATE en transacciones críticas
    \end{itemize}
\end{enumerate}

\subsection{Supuestos de Seguridad}

\begin{itemize}[leftmargin=*]
    \item Las claves privadas del servidor están protegidas mediante variables de entorno
    \item El personal de Servicio Social verifica identidad física durante check-in
    \item La comunicación cliente-servidor usa HTTPS (TLS 1.3)
    \item La base de datos está protegida por firewall y acceso restringido
    \item Los dispositivos de los socioformadores están bajo su control físico
\end{itemize}

\section{Criterios de Éxito}

\subsection{Funcionales}

\begin{itemize}[leftmargin=*]
    \item El sistema previene exitosamente inscripciones remotas (sin check-in)
    \item Los estudiantes no pueden inscribirse en más de un proyecto por periodo
    \item Los límites de capacidad se respetan en todos los escenarios
    \item Los códigos expiran correctamente y son de un solo uso
    \item Los recibos firmados son verificables y detectan manipulación
    \item Las exportaciones contienen datos completos y correctos
\end{itemize}

\subsection{No Funcionales}

\begin{itemize}[leftmargin=*]
    \item \textbf{Performance:} Carga de página < 2s, redención de código < 500ms
    \item \textbf{Disponibilidad:} Uptime > 99\% durante periodo de feria activa
    \item \textbf{Usabilidad:} Socioformadores pueden generar códigos sin capacitación
    \item \textbf{Seguridad:} Cero vulnerabilidades críticas en pruebas de seguridad
    \item \textbf{Escalabilidad:} Soporte para 500+ estudiantes simultáneos
\end{itemize}

\subsection{Académicos}

\begin{itemize}[leftmargin=*]
    \item Demostración clara del componente criptográfico
    \item Documentación técnica completa de arquitectura
    \item Código fuente bien estructurado y comentado
    \item Pruebas automatizadas con cobertura > 70\%
    \item Presentación efectiva con demo en vivo funcional
\end{itemize}

\section{Conclusión}

Este proyecto aborda un problema real del departamento de Servicio Social mediante una solución técnicamente sólida que integra conceptos de seguridad, criptografía, arquitectura de sistemas y desarrollo full-stack. El enfoque en verificación física mediante check-in, combinado con códigos efímeros y recibos firmados digitalmente, proporciona un sistema robusto que cumple tanto los requisitos operativos como los académicos del curso de criptografía.

\end{document}
